\documentclass{guide-ethmun}

\usepackage[ampersand]{easylist}
\usepackage{helvet}
\usepackage{marvosym}
\usepackage{hyperref}
\usepackage{dashrule}
\usepackage{pdfpages}
\usepackage{marvosym}
\usepackage{textcomp}

\newcommand{\quelle}[1]{\textit{\href{#1}{#1}}} % Hyperlink for sources, see References

\newcommand{\eur}[1]{\EUR\,#1\,bln.}            % € #1 bln.
\newcommand{\eurm}[1]{\EUR\,#1\,mln.}           % € #1 mln.
\newcommand{\usd}[1]{\textdollar\,#1\,bln.}     % $ #1 bln.
\newcommand{\usdm}[1]{\textdollar\,#1\,mln.}    % $ #1 mln.


\begin{document}

%%%%%%%%%%%%%%%%%%%%%%%%%%%%%%%%%%%%%%%%%%%%%%%%%%%%%%%%%%%%%%%%%%%

% Fill in Title Info here: 

\supertitle{Study Guide <semester>: <committee>}

\title{<committee>: <specific country>}

\author{\au{<name of author>}}

\address{\email{<email of author>}}

\date{logo-un.png} % File name of the Committee logo ("logo-un.png" for default UN logo, "ethmun-logo.png" for default ETHMUN logo)

\CYear{2020} % Adjust the correct Year

%%%%%%%%%%%%%%%%%%%%%%%%%%%%%%%%%%%%%%%%%%%%%%%%%%%%%%%%%%%%%%%%%%%

% \begin{abstract}
% We don't do abstracts at the moment.
% \end{abstract}

\maketitle

%%%%%%%%%%%%%%%%%%%%%%%%%%%%%%%%%%%%%%%%%%%%%%%%%%%%%%%%%%%%%%%%%%%

% Writing starts here: 



\section{Introduction and Executive Summary}\label{sec:summary}


This is an introduction \cite{ref:consilium}. 
\newline

\noindent This is a list of items:
\vspace{0.5em}
\begin{easylist}[itemize]
    & first item,
    & second item,
    & and many more...
\end{easylist}
\vspace{0.5em}

\noindent This paragraph has no indent. 

\section{Societal and Economic History}

\begin{table}[b]
\processtable{This is the heading of the table \cite{ref:table} \label{tab:table}}
{\begin{tabular*}{20pc}{@{\extracolsep{\fill}}llrr@{}}\toprule
Column 1 & Column 2 & Column 3 & Column 4\\
\midrule
Element 11  & Element 12 & Element 13 & Element 14\\
Element 21  & Element 22 & Element 23 & Element 24\\
Element 31  & Element 32 & Element 33 & Element 34\\
Element 41  & Element 42 & Element 43 & Element 44\\
\botrule
\end{tabular*}}{}
\end{table}

\section{Political Position}

\section{Foreign Policy \cite{ref:consilium}}

\newpage

\section{Domestic Policy}

\begin{figure}[!t]
\centering{\includegraphics[width=\columnwidth]{sample_map.png}}
\caption{A sample figure.}\label{fig:pipelines}
\end{figure}

\newpage

\section{Recent Developments}\label{recent}


%\begin{figure}[!t]
%\centering{\includegraphics{sample_1.eps}}
%\caption{Sample graph with blue (dotted), green (solid) and red (dashed) lines
%\figfooter{a}{Subfigure 1}
%\figfooter{b}{Subfigure 2}\label{fig1}}
%\end{figure}

%\begin{figure*}[!b]
%\centering{\includegraphics{sample_2.eps}}
%\caption{Sample graph spanning two columns\label{fig2}}
%\end{figure*}




%%%%%%%%%%%%%%%%%%%%%%%%%%%%%%%%%%%%%%%%%%%%%%%%%%%%%%%%%%%%%%%%%%%

% References start here:

% \vfill\pagebreak

\begin{thebibliography}{\hspace{1cm}}
\bibitem{ref:consilium}
European Council: 'Website of the European Council' \\
\quelle{https://www.consilium.europa.eu/en/}
\newline

\bibitem{ref:table}
<Author>: 'Title of the Reference' \\
\quelle{https://www.link-to-the-source.com/whatever/something/table}

\end{thebibliography}

\end{document}
